% from ega0-22-fr.tex
\oldpage[0\textsubscript{IV-1}]{211}

par (16.5.11). Or, comme $A_{\mathfrak p}$ est régulier (17.3.2),
$\dim(A_{\mathfrak p})$ est égal au rang sur $L$ de
$\mathfrak pA_{\mathfrak p}/\mathfrak p^2A_{\mathfrak p}$ (17.1.1), donc au
rang sur $L$ de $(\mathfrak p/\mathfrak p^2)\otimes_AL$, ce qui achève de
prouver que $i$ est injectif.
\end{proof}

\begin{theorem}[22.7.3]
\label{0.22.7.3-fr}
Soient $k$ un corps, $A$ un anneau local noethérien complet de corps résiduel
$K$. On suppose que~:
\begin{enumerate}
  \item[$1^\circ$] $[k:k^p]<+\infty$ (où $p$ est l'exposant caractéristique
  de $k$)~;
  \item[$2^\circ$] $K$ est une extension finie d'une extension séparable
  $K_0$ de $k$~;
  \item[$3^\circ$] $A$ est muni d'une structure de $K_0$-algèbre formellement
  lisse (pour la topologie préadique).
\end{enumerate}
Soient $\mathfrak q$ un idéal de $A$, $B=A/\mathfrak q$, $\mathfrak p$ un
idéal premier de $A$ contenant $\mathfrak q$. Les conditions suivantes sont
équivalentes~:
\begin{enumerate}
  \item[a)] $B_{\mathfrak p}$ est une $k$-algèbre formellement lisse (pour la
  topologie $\mathfrak p$-préadique).
  \item[b)] Il existe des $k$-dérivations $D_i$ de $A$ dans lui-même
  $(1\leq i\leq m)$ et des éléments $f_i$ $(1\leq i\leq m)$ de
  $\mathfrak q$, tels que les images des $f_i$ dans
  $\mathfrak qA_{\mathfrak p}$ engendrent cet idéal de $A_{\mathfrak p}$, et
  que l'on ait $\det(D_i(f_j))\notin\mathfrak p$.
  \item[b')] Il existe une sous-extension $k'$ de $K_0$, contenant $K_0^p$,
  telle que $[K_0:k']<+\infty$, des $k'$-dérivations $D_i$ de $A$ dans
  lui-même $(1\leq i\leq m)$ et des éléments $f_i$ de $\mathfrak q$
  $(1\leq i\leq m)$, tels que les images des $f_i$ dans
  $\mathfrak qA_{\mathfrak p}$ engendrent cet idéal de $A_{\mathfrak p}$, et
  que l'on ait $\det(D_i(f_j))\notin\mathfrak p$.
\end{enumerate}
\end{theorem}

\begin{proof}
Distinguons deux cas, suivant que $p=1$ ou $p>1$.

\emph{A)} $k$ \emph{est de caractéristique} $0$. Comme $K$ est alors une
extension séparable de $K_0$, il résulte de (19.6.4) que $A$ est
$K_0$-isomorphe à un anneau de séries formelles $K[[T_1,\ldots,T_r]]$ muni de
sa structure usuelle de $K$-algèbre. Mais alors, compte tenu de (19.8.8, (ii))
appliqué à $A/\mathfrak p$, on voit que les conditions générales de (22.7.2)
sont vérifiées en y remplaçant $k$ par $K$. En outre, en vertu de (19.6.4), il
revient au même de dire que $B_{\mathfrak p}$ est une $k$-algèbre
formellement lisse ou une $K$-algèbre formellement lisse, les deux conditions
étant équivalentes au fait que $B_{\mathfrak p}$ est un anneau régulier. On
peut donc appliquer les conclusions de (22.7.2), et il est immédiat que cela
prouve l'équivalence de a), b) et b') (avec $k'=K_0$ dans b')).

\emph{B)} $k$ \emph{est de caractéristique} $p>0$. Comme $A$ est une
$K_0$-algèbre formellement lisse et que $K_0$ est séparable sur $k$, il
résulte de (19.3.5, (ii)) et (19.6.1) que $A$ est une $k$-algèbre
formellement lisse~; en vertu de (22.5.9), $A_{\mathfrak p}$ est aussi une
$k$-algèbre formellement lisse pour la topologie $\mathfrak p$-préadique. Il
résulte donc encore de (22.6.2, (ii)) que la condition a) équivaut au fait que
l'homomorphisme $j_0$ défini dans (22.7.2.1) (où $L=k(\mathfrak p)$) est
injectif. Donc (19.1.12, c)) montre donc déjà que b) implique a)~; comme par
ailleurs b') implique trivialement b), tout revient à montrer que l'hypothèse
que $j_0$ est injectif entraîne b').

Désignons d'abord par $k'$ une sous-extension quelconque de $K_0$ telle que
$[K_0:k']<+\infty$. Notons que $\widehat\Omega^1_{A/k'}$ est un $A$-module
libre de type fini. Par le raisonnement de (22.7.2), il suffit de montrer que
(pour un choix convenable de $k'$), l'homomorphisme composé
\begin{equation}
\label{0.22.7.3.1-fr}
  j':(\mathfrak q/\mathfrak q^2)\otimes_AL
  \xrightarrow{j_0}\Omega^1_{A/k}\otimes_AL
  \longrightarrow\widehat\Omega^1_{A/k'}\otimes_AL
  \tag{22.7.3.1}
\end{equation}

